%=============================================================================!
% 7.6  RUNNING THE MOTION BACKWARDS                                           !
%=============================================================================!
\section{Running the motion backwards}
\label{sec:ha-reversal}

A film of a swinging pendulum looks as natural when it is run
backwards; a film of a book sliding to rest across a table does not,
because books are never seen to gather speed from a table that has cooled
a little. This section shows that the difference lies in the drag, and
that motion under a time-independent Hamiltonian can be reversed when
reversing every momentum leaves that Hamiltonian unchanged.

\subsection*{The reversed motion}

Let $q(t)$ and $p(t)$ be a motion. The reversed motion visits the same
positions in the opposite order, with every momentum reversed:
\begin{equation*}
   Q(t) = q(-t) , \qquad P(t) = -p(-t) .
\end{equation*}
Suppose that $H$ does not contain $t$ and that $H(q, -p) = H(q, p)$, as
it is whenever $K$ depends on the momenta through their squares. Differentiating this identity with respect
to $p$, by the chain rule on the left, gives $-(\partial H/\partial p)(q,
-p) = (\partial H/\partial p)(q, p)$, so reversing the momentum reverses
the sign of $\partial H/\partial p$; differentiating it with respect to
$q$ gives $(\partial H/\partial q)(q, -p) = (\partial H/\partial q)(q,
p)$, so reversing the momentum leaves $\partial H/\partial q$ unchanged.
Then, by the chain rule with the inner function $-t$, whose rate is $-1$,
and Hamilton's equations for the original motion at the time $-t$,
\begin{align*}
   \frac{\dd Q}{\dd t} &= -\dot{q}(-t)
   = -\frac{\partial H}{\partial p}\bigl(q(-t), p(-t)\bigr)
   = \frac{\partial H}{\partial p}(Q, P) ,\\
   \frac{\dd P}{\dd t} &= \dot{p}(-t)
   = -\frac{\partial H}{\partial q}\bigl(q(-t), p(-t)\bigr)
   = -\frac{\partial H}{\partial q}(Q, P) ,
\end{align*}
where the last step of each line uses $Q = q(-t)$, $P = -p(-t)$ and the two
results just found. The reversed motion obeys Hamilton's equations too.
Since $H$ does not contain $t$, shifting the time changes nothing in these
equations, so $q(\tau - t)$, $-p(\tau - t)$ obeys them as well. It starts
at $\bigl(q(\tau), -p(\tau)\bigr)$, and a state fixes the motion that
follows it (\cref{sec:ne-equations}), so if a motion is stopped at a time
$\tau$, every momentum is reversed, and the motion is followed for another time $\tau$, the system retraces
its path and arrives where it started, with its starting momenta
reversed.

\Cref{fig:ha-reversal}(a) shows this for a puck of \SI{0.2}{\kilogram}
on smooth ice, held near a point by springs that are stiffer north--south
(\SI{5}{\newton\per\metre}) than east--west (\SI{2}{\newton\per\metre}),
so that it traces a looping path. Followed for \SI{4}{\second}, reversed
and followed for another \SI{4}{\second}, it returns to within $3\times
10^{-14}$~\si{\metre} of its start, with the momentum
$(-0.02, -0.10)$~\si{\kilogram\metre\per\second}, the reverse of the one it
started with.

\begin{figure}[htb]
\centering
\includegraphics{ch07_hamilton/reversal}
\caption{A puck of \SI{0.2}{\kilogram} on springs of 2 (east--west) and
\SI{5}{\newton\per\metre} (north--south), started at A with the velocity
$(0.1, 0.5)$~\si{\metre\per\second}, followed for \SI{4}{\second} to B
(teal), then with every momentum reversed for another \SI{4}{\second}
(ochre, dashed), ending at the square. (a) Without drag it retraces its
path to A. (b) With drag at $\gamma = \SI{0.3}{\per\second}$ it does
not.}
\label{fig:ha-reversal}
\end{figure}

\subsection*{Drag breaks the symmetry}

With drag, $\dot{p} = -\partial H/\partial q - \gamma p$. Repeating the
second line above, the reversed motion has $\dd P/\dd t = \dot{p}(-t) =
-\partial H/\partial q - \gamma p(-t) = -\partial H/\partial q + \gamma P$:
it would need a force that pushes the puck along its motion, in place of
one that resists it, and the dragged equations do not supply one.
\Cref{fig:ha-reversal}(b) shows the result: after the reversal the puck
keeps losing energy and ends \SI{0.21}{\metre} from its start. Sliding friction,
like drag, always opposes the motion, so the reversed slide would need it
to push instead; it is what makes the book's slide irreversible, and the
motion of atoms under forces from a potential has no such term.

Chapter~9 asks the same of a method of computing motion: one step
forwards, a reversal of the momenta and one more step should bring the
system back to where it began, as the exact motion does.

\begin{keyidea}
When $H$ is time-independent and unchanged by reversing the momenta, the reversed motion $q(-t)$,
$-p(-t)$ obeys Hamilton's equations, so reversing every momentum makes a
system retrace its path. Drag breaks this.
\end{keyidea}

\begin{selfcheck}
A ball thrown straight up at \SI{12}{\metre\per\second} reaches its highest
point after \SI{1.224}{\second}. Using time reversal, and without solving
any equation, say when and how fast it comes back to the hand.
\end{selfcheck}
